\documentclass[10pt]{article}
\usepackage[margin=20mm]{geometry}
\usepackage{amsmath,amssymb,mathtools,fontspec,unicode-math,xurl,hyperref}
\IfFontExistsTF{DejaVu Serif}{\setmainfont{DejaVu Serif}}{\setmainfont{Latin Modern Roman}}
\setmathfont{latinmodern-math.otf}
\setlength{\parindent}{0pt}\setlength{\parskip}{6pt}\sloppy
\title{Initial scientific review
A half-space Bernstein theorem
for anisotropic minimal graphs}
\author{AI-assisted technical review}
\date{}
\begin{document}
\maketitle
Assignment. Authors. Wenkui Du, Connor Mooney, Yang Yang and Jingze Zhu. Exact source. arXiv:2312.07519v1, 12 December 2023, 18:51:13 UTC; nine pages. The journal record is 10.4171/JEMS/1695 (https://doi.org/10.4171/JEMS/1695), published online on 29 August 2025. The frozen preprint is not assumed equivalent to the later publication.\par
\section*{Assessment and coverage}
No Issue Found within the declared review and focused acceptance scope. No substantive mathematical defect was established. The body–boundary notation in (2) is clarified below, but the subsequent proof uses the intended boundary identity. Since no essential estimate is invalidated, this clarification is not counted as a numbered Problem. In particular, the apparent nonsymmetry problem in the contact-set Jacobian does not furnish a counterexample to the determinant estimate: the appropriate similarity transformation proves that estimate.\par
All nine pages, including references and all six figures, were read and visually inspected. The review followed the support-function geometry, graph calibration, radial barrier, Wulff-contact area estimate, slope-improvement lemma, the three domain cases and the two infinite-slope subcases. The hypotheses \(u\in C^\infty(\Omega)\cap C(\overline\Omega)\), convex proper domain, linear boundary data and smooth uniformly convex Wulff shape were retained. No growth bound at infinity or symmetry of \(\Phi\) was silently added.\par
\section*{A notation clarification}
Equation (2): the Wulff body and its boundary. On p. 2, (2) prints \(K=\nabla\Phi(S^n)\). On p. 1, \(K\) is the convex Wulff body, with the area case explicitly given by \(K=B_1\). For the admissible support function \(\Phi(z)=|z|\), \[\nabla\Phi(x)=x\quad(x\in S^n),\qquad
 \nabla\Phi(S^n)=S^n,\] whereas \(0\in B_1\) and \(0\notin S^n\). Thus the displayed identity is false for the stated body notation. Its intended replacement is \[\partial K=\nabla\Phi(S^n).\] The inverse Gauss map and the curvature identities concern this boundary. Subsequent uses of translated/rescaled Wulff shapes have the usual boundary interpretation, so no failure of the measure estimate or of Theorem 1.1 follows from this correction. The evidence is an exact admissible example, not an interpretation based on extracted text alone: the original pages were inspected.\par
\section*{Independent checks of the central estimates}
\subsection*{The contact Jacobian}
Let \(A=D_T^2\Phi(\nu)>0\) and \(B=II_S\), using the sign convention of the contact inequality on pp. 3–4. The argument needs \[B\ge-r^{-1}A^{-1},\qquad \operatorname{tr}(AB)=0,\qquad
 J=I+rAB.\] Although \(J\) need not be symmetric in the Euclidean metric, \[A^{-1/2}JA^{1/2}=I+rA^{1/2}BA^{1/2}\] is symmetric positive semidefinite with trace \(n\). Its eigenvalues are therefore nonnegative and have sum \(n\), giving \[0\le\det J\le1.\] Equivalently, \(J\) is self-adjoint and nonnegative in the \(A^{-1}\) inner product. This supplies the intended meaning of the shorthand \(D_xy\ge0\).\par
As an exact stress test, take \[A=\begin{pmatrix}16&0\\0&1\end{pmatrix},\quad
 B=\begin{pmatrix}0&1/8\\1/8&0\end{pmatrix},\quad r=1.\] Then \(B+A^{-1}>0\), with determinant \(3/64\), while \[J=\begin{pmatrix}1&2\\1/8&1\end{pmatrix},\qquad
 (1,-1)J(1,-1)^T=-1/8.\] Nevertheless its eigenvalues are \(1/2,3/2\), and \(\det J=3/4\). These rational calculations were independently executed. They refute a Euclidean quadratic-form reading of the shorthand, while confirming the actual determinant bound. This distinction is why no additional substantive finding is recorded here.\par
For the geometric map \(y=x-r\nabla\Phi(\nu(x))\), differentiation gives \(Dy=I-rA\,D\nu\). With \(II=-D\nu\) this is the printed plus sign; with \(II=D\nu\) the sign and contact inequality must both reverse. The proof can be expressed consistently with either convention. The manuscript would benefit from stating that convention explicitly and using it for both the Wulff boundary and \(S\); changing only one sign is not a counterexample to the estimate.\par
At nonsingular contact points the image differential lies in the same tangent plane \(\nu^\perp\). Hence projection in the vertical direction supplies the same cosine factor on both sides. Together with first-contact selection and the small normal tilt, the determinant bound has the required direction for the area formula.\par
\subsection*{Radial barriers and the scales}
For \(f(\rho)=\rho^{-M}\), the radial Hessian eigenvalue is \(M(M+1)\rho^{-M-2}\) and each of the \(n-1\) tangential eigenvalues is \(-M\rho^{-M-2}\). Thus the Pucci convention on p. 3 yields \[\mathcal M^-_\lambda(D^2f)
  =M\rho^{-M-2}
       \bigl(\lambda(M+1)-(n-1)/\lambda\bigr)>0\] once \(M+1>(n-1)/\lambda^2\). For fixed \(\delta\), taking \(\varepsilon M\delta^{-M-1}<1\) keeps the gradient within the ellipticity region. The truncation at \(\rho=\delta\) is harmless in the proposed contradiction, since the excluded close-contact region is precisely the inner cylinder.\par
The Wulff scale \(r\asymp\delta^{1/2}\) gives a height variation \(\delta^2/r\asymp\delta^{3/2}\) and normal variation \(\delta/r\asymp\delta^{1/2}\). The lost horizontal margin is \(O(r)\). These are the three exponents in Lemma 2.1, and they are mutually consistent. The choice order is fixed \(\delta\), then sufficiently small \(\varepsilon(\delta,n,K)\); the argument does not require an unjustified uniform choice as \(\delta\to0\).\par
\subsection*{Calibration and thin-cylinder competitors}
For a graph, homogeneity gives \(\nabla\Phi(\nu)=\nabla\Phi(-\nabla w,1)\). Its first \(n\) components are \(-\nabla\varphi(\nabla w)\). Extending the field vertically therefore gives \[\operatorname{div}\nabla\Phi(\nu)
     =-\operatorname{div}_x\nabla\varphi(\nabla w)=0.\] The support inequality \(V\cdot a\le\Phi(a)\) and equality at the graph normal give the claimed calibration comparison for compact changes away from the vertical sides. This restriction is material. In the slab case the upper comparison cylinder avoids those sides for large \(k\), because the prescribed boundary heights tend to those of \(\Gamma\).\par
If positivity of the integrand is used in the extra-sheet lower bound, one may translate the Wulff body so that the origin lies in its interior. Replacing \(\Phi(\nu)\) by \(\Phi(\nu)-q\cdot\nu\) changes neither the Euler equation nor energy differences for compact competitors, by the divergence theorem. Thus this normalization is available without assuming that the original body was centered.\par
The competitor costs in pp. 6–8 have the correct distinction. In a shrinking slab there is no far \(x_1\)-face cost. In the half-space there is at most one face, with normal \(-e_1\), in addition to the \(O(h)\) thin sides. The latter face cancels the leading sheet contribution, so the extra projected sheet is essential. Reflection also changes the anisotropy as stated in the paper; no evenness assumption on \(\Phi\) is needed.\par
\section*{Supporting checks}
The definitions of \(A_\pm\) allow the stated extended-real possibilities because evaluation at any interior point bounds the opposite end. In the slab and the third convex-domain case, a boundary point with \(x_1>0\) enforces \(A_-\le0\le A_+\). Convexity and \(0\in\overline\Omega\) give \(k^{-1}\overline\Omega\subset\overline\Omega\); the closure bars lost in some text extraction are visible on the original pages.\par
In Lemma 3.1, a subsequence avoiding a fixed off-edge point supplies a fixed barrier improvement, so rescaling expands \(B_\delta\) to \(B_{k_j\delta}\). This contradicts the extremal slope definition. The requirement of a nonzero slope in cases (B) and (C) is used to keep the barriers below the prescribed boundary. Finite unequal half-space slopes give incompatible sheets after projection. For two infinite slopes, continuity along graphical columns forces an additional crossing in at least one of the upper and lower cylinders; reflection puts the useful competitor above the boundary edge.\par
These checks support the proof’s mechanism, but they are not a complete formal reconstruction of every geometric limiting argument. In particular, the precise first-contact/properness formulation of the local surface lemma, all quantitative boundary-barrier placements in Lemma 3.1, and the full measure-theoretic extraction of the extra sheet from the final pigeonhole argument were not independently rewritten as standalone theorems. No counterexample satisfying their intended geometric hypotheses was established. The nonconvex extension in Remark 3.2 was read, but not independently proved in its full generality. These are review limits, not confirmed failures.\par
Searches covered the exact arXiv version, the publisher and author research listings. No separate supplement or author code was located. The attempted current author-hosted PDF returned a 502 fetch error; another author’s page redirected to an inaccessible sign-in page. Neither was treated as read or as evidence of equivalence to the frozen source. No author code was executed.\par
The exact source, own calculations, TeX, existing compiler, dependency payload, PDF and page inspections are hash-bound in the accompanying records. Focused independent acceptance revisited the sole reported issue, its actual impact and the coverage boundaries, and removed it from substantive Problems. This is not a full second reconstruction of the geometric proof. Scientific adjudication remains pending, and human reviews are empty.\par
\end{document}
